% --- Archon physics formalization source begin ---
% archon:physics
% archon:covers PhyXMiniProblems/problem_phyx_mini_0676.lean
% archon:source-report reports/phyx_mini/problem_phyx_mini_0676.source.json

\chapter{Physics problem phyx\_mini\_0676}
\label{ch:PhyXMiniProblems_problem_phyx_mini_0676}

\paragraph{Problem source.}
\#\# Physical scenario

A surveyor measures the distance across a straight river by the following method as shown in figure. Starting directly across from a tree on the opposite bank, she walks \textbackslash{}( d = 100 \textbackslash{}text\{ m\} \textbackslash{}) along the riverbank to establish a baseline. Then she sights across to the tree. The angle from her baseline to the tree is \textbackslash{}( \textbackslash{}theta = 35.0\textasciicircum{}\textbackslash{}circ \textbackslash{}).

\#\# Question

How wide is the river?

\#\# Answer choices

- A: A: 24.0m

- B: B: 35.0m

- C: C: 70.0m

- D: D: 65.0m

\#\# Figure caption (auxiliary; use the image as primary evidence)

The image depicts a diagram of a river or a similar body of water. The main elements are:

1. **River**: A horizontal blue area representing the waterway.
2. **Land**: Green areas on both sides of the river, representing the riverbanks.
3. **Bush**: Located on the left riverbank.
4. **Points**: Two black dots indicating specific positions: one on the bush and one on the right riverbank.
5. **Line and Angle**: A diagonal line connects the two points across the river, with an angle \textbackslash{}( \textbackslash{}theta \textbackslash{}) formed between this line and a horizontal line on the right bank.
6. **Distance**: The horizontal line on the right riverbank is labeled with \textbackslash{}( d \textbackslash{}), indicating the distance between two vertical dashed lines that extend from the two points to the riverbank.

Overall, the diagram illustrates the geometric relationship and distance across the river.

\#\# Dataset metadata

Kinematics; Spatial Relation Reasoning

\paragraph{Recorded answer/context.}
C: C: 70.0m

\paragraph{Figure/image path.}
/root/proposal\_for\_physic/science-mango/phyx\_mini\_run/phyx\_data/test\_image/676.png

\paragraph{Formalization target.}
create a compiling Lean file with sorry bodies at `PhyXMiniProblems/problem\_phyx\_mini\_0676.lean`. The Lean declarations must preserve the physical quantities, dimensions or dimensional roles, figure labels, governing-law hypotheses, and final relation expressed by this problem.
Use Mathlib/Physlib names found through LeanExplore where available. If a physics API is missing, introduce faithful local abstractions rather than scalar placeholder aliases.

\begin{theorem}[Physics formalization target]
\label{thm:physics:phyx_mini_0676:target}
\lean{PhyXMiniProblems.ProblemPhyXMini0676.riverWidth_is_70_metres_answer_C}
\uses{def:physics:phyx-mini-0676:phyxminiproblems-problemphyxmini0676-lengthinmeters, def:physics:phyx-mini-0676:phyxminiproblems-problemphyxmini0676-riversurveysetup, def:physics:phyx-mini-0676:phyxminiproblems-problemphyxmini0676-matchesprimaryriversurveyfigure, def:physics:phyx-mini-0676:phyxminiproblems-problemphyxmini0676-hasphysicalriversurveyconfiguration, def:physics:phyx-mini-0676:phyxminiproblems-problemphyxmini0676-matchesproblemmeasurements, def:physics:phyx-mini-0676:phyxminiproblems-problemphyxmini0676-satisfiesrighttriangletangentlaw, def:physics:phyx-mini-0676:phyxminiproblems-problemphyxmini0676-matchesanswertonearesttenth, def:physics:phyx-mini-0676:phyxminiproblems-problemphyxmini0676-isuniqueclosestanswer}
The assigned autoformalize agent should translate this physics problem into Lean declarations in the covered file, with theorem and lemma proof bodies written as `by sorry`.
\end{theorem}
\begin{proof}
This is an autoformalization task, not a proof task. Produce statements that can later be proved by the physics prover without weakening the physical model.
\end{proof}

% --- Archon declaration topology begin ---
\section{Declaration topology}
\begin{definition}[Length Quantity]
  \label{def:physics:phyx-mini-0676:phyxminiproblems-problemphyxmini0676-lengthquantity}
  \lean{PhyXMiniProblems.ProblemPhyXMini0676.LengthQuantity}
  A physical length, independent of the unit system used to read it
\end{definition}
\begin{proof}
  This is the declared model definition of Length Quantity.
\end{proof}
\begin{definition}[length In Meters]
  \label{def:physics:phyx-mini-0676:phyxminiproblems-problemphyxmini0676-lengthinmeters}
  \lean{PhyXMiniProblems.ProblemPhyXMini0676.lengthInMeters}
  \uses{def:physics:phyx-mini-0676:phyxminiproblems-problemphyxmini0676-lengthquantity}
  The scalar SI readout of a physical length, in metres
\end{definition}
\begin{proof}
  This is the declared model definition of length In Meters.
\end{proof}
\begin{definition}[Survey Point]
  \label{def:physics:phyx-mini-0676:phyxminiproblems-problemphyxmini0676-surveypoint}
  \lean{PhyXMiniProblems.ProblemPhyXMini0676.SurveyPoint}
  The three labeled points in the river-survey construction
\end{definition}
\begin{proof}
  This is the declared model definition of Survey Point.
\end{proof}
\begin{definition}[River Bank Side]
  \label{def:physics:phyx-mini-0676:phyxminiproblems-problemphyxmini0676-riverbankside}
  \lean{PhyXMiniProblems.ProblemPhyXMini0676.RiverBankSide}
  The two sides of the river shown in the primary figure
\end{definition}
\begin{proof}
  This is the declared model definition of River Bank Side.
\end{proof}
\begin{definition}[River Survey Setup]
  \label{def:physics:phyx-mini-0676:phyxminiproblems-problemphyxmini0676-riversurveysetup}
  \lean{PhyXMiniProblems.ProblemPhyXMini0676.RiverSurveySetup}
  \uses{def:physics:phyx-mini-0676:phyxminiproblems-problemphyxmini0676-lengthquantity, def:physics:phyx-mini-0676:phyxminiproblems-problemphyxmini0676-surveypoint, def:physics:phyx-mini-0676:phyxminiproblems-problemphyxmini0676-riverbankside}
  The physical quantities and the smallest relational geometry interface needed to retain the scenario and primary-figure labels. baselineDistanceD is the walked bank distance labeled d, riverWidth is the perpendicular distance between the straight banks, and sightAngleThetaRadians is the angle labeled θ at the sighting station. The relational fields distinguish the figure's baseline and sight line from mere numerical lengths
\end{definition}
\begin{proof}
  This is the declared model definition of River Survey Setup.
\end{proof}
\begin{definition}[Matches Primary River Survey Figure]
  \label{def:physics:phyx-mini-0676:phyxminiproblems-problemphyxmini0676-matchesprimaryriversurveyfigure}
  \lean{PhyXMiniProblems.ProblemPhyXMini0676.MatchesPrimaryRiverSurveyFigure}
  \uses{def:physics:phyx-mini-0676:phyxminiproblems-problemphyxmini0676-riversurveysetup}
  The categorical and incidence information read from the primary diagram and the statement that the surveyor starts directly opposite the tree. No metric value of the river width occurs here
\end{definition}
\begin{proof}
  This is the declared model definition of Matches Primary River Survey Figure.
\end{proof}
\begin{definition}[Has Physical River Survey Configuration]
  \label{def:physics:phyx-mini-0676:phyxminiproblems-problemphyxmini0676-hasphysicalriversurveyconfiguration}
  \lean{PhyXMiniProblems.ProblemPhyXMini0676.HasPhysicalRiverSurveyConfiguration}
  \uses{def:physics:phyx-mini-0676:phyxminiproblems-problemphyxmini0676-lengthinmeters, def:physics:phyx-mini-0676:phyxminiproblems-problemphyxmini0676-riversurveysetup}
  Positive physical lengths and the acute branch of the depicted angle
\end{definition}
\begin{proof}
  This is the declared model definition of Has Physical River Survey Configuration.
\end{proof}
\begin{definition}[Matches Problem Measurements]
  \label{def:physics:phyx-mini-0676:phyxminiproblems-problemphyxmini0676-matchesproblemmeasurements}
  \lean{PhyXMiniProblems.ProblemPhyXMini0676.MatchesProblemMeasurements}
  \uses{def:physics:phyx-mini-0676:phyxminiproblems-problemphyxmini0676-lengthinmeters, def:physics:phyx-mini-0676:phyxminiproblems-problemphyxmini0676-riversurveysetup}
  The two calibrated measurements stated in the problem: a 100 m baseline and a 35° sight angle. The angle equality explicitly converts degrees to the radian readout used by Real.tan
\end{definition}
\begin{proof}
  This is the declared model definition of Matches Problem Measurements.
\end{proof}
\begin{definition}[Satisfies Right Triangle Tangent Law]
  \label{def:physics:phyx-mini-0676:phyxminiproblems-problemphyxmini0676-satisfiesrighttriangletangentlaw}
  \lean{PhyXMiniProblems.ProblemPhyXMini0676.SatisfiesRightTriangleTangentLaw}
  \uses{def:physics:phyx-mini-0676:phyxminiproblems-problemphyxmini0676-lengthinmeters, def:physics:phyx-mini-0676:phyxminiproblems-problemphyxmini0676-riversurveysetup}
  The governing right-triangle relation tan θ = width / d, written without a division by the positive baseline length. This is the generic triangulation law, not the requested numerical conclusion
\end{definition}
\begin{proof}
  This is the declared model definition of Satisfies Right Triangle Tangent Law.
\end{proof}
\begin{definition}[Answer Choice]
  \label{def:physics:phyx-mini-0676:phyxminiproblems-problemphyxmini0676-answerchoice}
  \lean{PhyXMiniProblems.ProblemPhyXMini0676.AnswerChoice}
  Labels of the four multiple-choice answers in the problem statement
\end{definition}
\begin{proof}
  This is the declared model definition of Answer Choice.
\end{proof}
\begin{definition}[answer Width Meters]
  \label{def:physics:phyx-mini-0676:phyxminiproblems-problemphyxmini0676-answerwidthmeters}
  \lean{PhyXMiniProblems.ProblemPhyXMini0676.answerWidthMeters}
  \uses{def:physics:phyx-mini-0676:phyxminiproblems-problemphyxmini0676-answerchoice}
  The river-width readout, in metres, printed beside each answer choice
\end{definition}
\begin{proof}
  This is the declared model definition of answer Width Meters.
\end{proof}
\begin{definition}[Matches Answer To Nearest Tenth]
  \label{def:physics:phyx-mini-0676:phyxminiproblems-problemphyxmini0676-matchesanswertonearesttenth}
  \lean{PhyXMiniProblems.ProblemPhyXMini0676.MatchesAnswerToNearestTenth}
  \uses{def:physics:phyx-mini-0676:phyxminiproblems-problemphyxmini0676-lengthquantity, def:physics:phyx-mini-0676:phyxminiproblems-problemphyxmini0676-lengthinmeters, def:physics:phyx-mini-0676:phyxminiproblems-problemphyxmini0676-answerchoice, def:physics:phyx-mini-0676:phyxminiproblems-problemphyxmini0676-answerwidthmeters}
  A physical width rounds to a displayed answer at one decimal place
\end{definition}
\begin{proof}
  This is the declared model definition of Matches Answer To Nearest Tenth.
\end{proof}
\begin{definition}[Is Unique Closest Answer]
  \label{def:physics:phyx-mini-0676:phyxminiproblems-problemphyxmini0676-isuniqueclosestanswer}
  \lean{PhyXMiniProblems.ProblemPhyXMini0676.IsUniqueClosestAnswer}
  \uses{def:physics:phyx-mini-0676:phyxminiproblems-problemphyxmini0676-answerchoice, def:physics:phyx-mini-0676:phyxminiproblems-problemphyxmini0676-answerwidthmeters}
  A choice is strictly closer to a metre readout than every other choice
\end{definition}
\begin{proof}
  This is the declared model definition of Is Unique Closest Answer.
\end{proof}
% --- Archon declaration topology end ---
% --- Archon physics formalization source end ---
